

\subsection*{\SGAInlineRefTarget{sga3:iii:subsection:1-3}1.3. Extensions of group laws.}
Consider the following situation: one has a morphism in
\(\widehat{\mathscr C}\),
\[
 (\dagger)\qquad p:X\longrightarrow Y,
\]
and a commutative \(\widehat{\mathscr C}\)-group \(M\) acting on
\(X\), such that \(X\) is formally principal homogeneous over \(Y\)
under \(M_Y\).

If \(g:Y\to Z\) is an arbitrary morphism in
\(\widehat{\mathscr C}\), then \(g\circ p:X\to Z\) is invariant
under \(M\): for every \(S\in\operatorname{Ob}\mathscr C\), the map
\((g\circ p)(S)\) is invariant under the action of \(M(S)\) on
\(X(S)\). Conversely, we shall assume that the following condition
holds for \(n=1,2,3,4\):

\medskip
\noindent\((+)_{n}\): \textit{Every morphism from \(X^n\) to \(M\)
that is invariant under the action of \(M^n\) on \(X^n\) factors
uniquely through \(p^n:X^n\to Y^n\), where the superscript \(n\)
denotes a Cartesian power.}

\medskip
\noindent\SGAInlineRefTarget{sga3:iii:lemma:1-3-1}\textbf{Lemma 1.3.1.}
\begin{enumerate}
 \item[(i)] \textit{If \(h\) is a morphism from \(Y\) to \(M\),
 the automorphism \(u_h\) of \(X\), defined set-theoretically by
 \(x\mapsto h(p(x))\cdot x\), preserves the fibers of \(p\) and
 commutes with the operations of \(M\) on \(X\).\footnote{Editor's
 note: here and below, ``commutes with \(p\) and the operations of
 \(M\)'' has been replaced by ``preserves the fibers of \(p\) and
 commutes with the operations of \(M\).'' The equalities that follow
 have also been added.} Thus, for every
 \(S\in\operatorname{Ob}\mathscr C\), \(x\in X(S)\), and
 \(m\in M(S)\), one has}
 \[
  p\bigl(h(p(x))\cdot x\bigr)=p(x),
  \qquad
  m\cdot h(p(x))\cdot x
   =h(p(m\cdot x))\cdot m\cdot x.
 \]

 \item[(ii)] \textit{This construction gives a bijective
 correspondence between morphisms from \(Y\) to \(M\) and
 automorphisms of \(X\) that preserve the fibers of \(p\) and commute
 with the operations of \(M\).}
\end{enumerate}

The first assertion is clear, since \(p(m\cdot x)=p(x)\) and \(M\)
is commutative. Conversely, an automorphism \(u\) of \(X\) preserving
the fibers of \(p\) can be written set-theoretically as
\(x\mapsto g(x)x\), where \(g\) is a certain morphism from \(X\) to
\(M\). If \(u\) commutes with the operations of \(M\), then \(g\) is
invariant under \(M\), and the
conclusion follows from condition \((+)_{1}\).

Suppose now, in addition, that a group law on \(Y\) and an action of
\(Y\) on \(M\) are given, that is, a morphism of
\(\widehat{\mathscr C}\)-groups
\[
 (\ddagger)\qquad
 f:Y\longrightarrow
 \underline{\operatorname{Aut}}_{\widehat{\mathscr C}\text{-gr.}}(M).
\]

\medskip
\noindent\SGAInlineRefTarget{sga3:iii:definition:1-3-2}\textbf{Definition 1.3.2.}
\textit{A composition law on \(X\),}
\[
 P:X\times X\longrightarrow X,
\]
\textit{is called \emph{admissible} if it satisfies the following two
conditions:}

\begin{enumerate}
 \item[\SGAInlineRefTarget{sga3:iii:definition:1-3-2:item:i}(i)] \textit{\(P\) lifts the group law of \(Y\); that is, the
 following diagram is commutative.}

 \SGAThreeNativeDiagram[0.28\textwidth]
  {native_diagrams/exp3/diagram_029_exp3_p029_admissible_law_square.tex}
  {The compatibility of an admissible composition law with the group
   law on \(Y\).}
  {fig:exp3-p029-admissible-law-square}
  {LNM 151, pp.~108--109}

 \item[\SGAInlineRefTarget{sga3:iii:definition:1-3-2:item:ii}(ii)] \textit{For every \(S\in\operatorname{Ob}\mathscr C\)
 and all \(x,y\in X(S)\), \(m,n\in M(S)\), one has}
 \[
  (++)\qquad
  P(m\cdot x,n\cdot y)
   =m\cdot f(p(x))(n)\cdot P(x,y).
 \]
\end{enumerate}

\medskip
\noindent\SGAInlineRefTarget{sga3:iii:proposition:1-3-3}\textbf{Proposition 1.3.3.}
\textit{For a group law \(\ast\) on \(X\) to be admissible, it is
necessary and sufficient that the following four conditions hold:}
\begin{enumerate}
 \item[(i)] \textit{\(p:X\to Y\) is a group morphism.}

 \item[(ii)] \textit{The morphism \(i:M\to X\), defined by
 \(i(m)=m\cdot e_X\), is a group isomorphism from \(M\) onto
 \(\ker(p)\); in other words, set-theoretically,}
 \[
  (m\cdot e_X)\ast(n\cdot e_X)=(mn)\cdot e_X.
 \]

 \item[(iii)] \textit{One has
 \(m\cdot x=(m\cdot e_X)\ast x=i(m)\ast x\) for every
 \(m\in M(S)\) and \(x\in X(S)\).}

 \item[(iv)] \textit{The inner automorphisms of \(X\) act on
 \(\ker(p)\) according to the set-theoretic formula}
 \[
  x\ast i(m)\ast x^{-1}=i\bigl(f(p(x))m\bigr).
 \]
\end{enumerate}
The proof is immediate.

\medskip
\noindent\SGAInlineRefTarget{sga3:iii:lemma:1-3-4}\textbf{Lemma 1.3.4.}\footnote{Editor's note: the statement
and its proof have been expanded with a view to step~\SGARefLink{sga3:iii:paragraph:2-0:equation:e}{(e)} in the proof
of \SGARefLink{sga3:iii:proposition:1-3-6}{1.3.6}.}
\textit{Let \(h:Y\to M\) be a morphism and let \(u_h\) be the
automorphism \(x\mapsto h(p(x))\cdot x\) of \(X\) (cf.~\SGARefLink{sga3:iii:lemma:1-3-1}{1.3.1}). Let
\(P\) be an admissible composition law (respectively, an admissible
group law) on \(X\), and let \(P'\) be the composition law on \(X\)
obtained from \(P\) by means of \(u_h\), that is,}
\[
 P'(x,y)=u_h^{-1}\bigl(P(u_h(x),u_h(y))\bigr).
\]
\textit{Then:}
\begin{enumerate}
 \item[\SGAInlineRefTarget{sga3:iii:lemma:1-3-4:item:i}(i)] \textit{\(P'\) is an admissible composition law
 (respectively, an admissible group law).}

 \item[\SGAInlineRefTarget{sga3:iii:lemma:1-3-4:item:ii}(ii)] \textit{For every \(x,y\in X(S)\), with
 \(S\in\operatorname{Ob}\mathscr C\), set \(v=p(x)\) and
 \(w=p(y)\). Then}
 \[
  P'(x,y)=h(vw)^{-1}\cdot h(v)\cdot f(v)(h(w))\cdot P(x,y)
   =(\partial^1h)(p(x),p(y))\cdot P(x,y).
 \]
\end{enumerate}

\noindent\textit{Proof.}
One has \(u_h^{-1}=u_{h^\vee}\), where \(h^\vee:Y\to M\) is
defined by \(h^\vee(y)=h(y)^{-1}\). By \SGARefLink{sga3:iii:definition:1-3-2:item:i}{1.3.2(i)} and \SGARefLink{sga3:iii:lemma:1-3-4:item:ii}{(ii)},
\[
 P(h(v)\cdot x,h(w)\cdot y)
  =h(v)\cdot f(v)(h(w))\cdot P(x,y),
 \qquad p(P(x,y))=vw,
\]
and hence
\[
 P'(x,y)=h(vw)^{-1}\cdot h(v)\cdot f(v)(h(w))\cdot P(x,y)
  =(\partial^1h)(p(x),p(y))\cdot P(x,y).
\]
It follows immediately that \(P'\) satisfies conditions \SGARefLink{sga3:iii:lemma:1-3-4:item:i}{(i)} and \SGARefLink{sga3:iii:lemma:1-3-4:item:ii}{(ii)}
of \SGARefLink{sga3:iii:definition:1-3-2}{1.3.2}.

\medskip
\noindent\SGAInlineRefTarget{sga3:iii:definition:1-3-5}\textbf{Definition 1.3.5.}
\textit{Two admissible composition laws obtained from one another by
the procedure of \SGARefLink{sga3:iii:lemma:1-3-4}{1.3.4} are called \emph{equivalent}.}

\medskip
\noindent\SGAInlineRefTarget{sga3:iii:proposition:1-3-6}\textbf{Proposition 1.3.6.}
\textit{Suppose that an admissible composition law on \(X\) exists.
Then:}
\begin{enumerate}
 \item[\SGAInlineRefTarget{sga3:iii:proposition:1-3-6:item:i}(i)] \textit{There exists a canonically determined class
 \(c\in H^3(Y,M)\), whose vanishing is necessary and sufficient for
 the existence of an admissible associative composition law on \(X\).}

 \item[\SGAInlineRefTarget{sga3:iii:proposition:1-3-6:item:ii}(ii)] \textit{If \(c=0\), the set of admissible associative
 composition laws on \(X\) (respectively, the set of their equivalence
 classes) is principal homogeneous under \(Z^2(Y,M)\) (respectively,
 \(H^2(Y,M)\)).}
\end{enumerate}

The proof proceeds in several steps.

\smallskip
\noindent\textbf{(a)}
Let \(P\) be an admissible composition law on \(X\). Since \(P\)
lifts the associative composition law of \(Y\), there is a unique
morphism \(a:X^3\to M\) such that
\[
 (*)\qquad
 P(x,P(y,z))=a(x,y,z)\cdot P(P(x,y),z).
\]
Applying conditions \SGARefLink{sga3:iii:definition:1-3-2:item:i}{1.3.2(i)} and \SGARefLink{sga3:iii:proposition:1-3-6:item:ii}{(ii)}, one sees at once that \(a\) is
invariant under the action of \(M^3\) on \(X^3\).\footnote{Editor's
note: for \(x,y,z\in X(S)\), set \(u=p(x)\), \(v=p(y)\), and
\(w=p(z)\). Then
\[
\begin{aligned}
 P(mx,P(m'y,m''z))
  &=P(mx,m'f(v)(m'')P(y,z))\\
  &=mf(u)(m')f(uv)(m'')P(x,P(y,z)).
\end{aligned}
\]
On the other hand, since \(p(P(x,y))=uv\), one also has
\[
\begin{aligned}
 P(P(mx,m'y),m''z)
  &=P(mf(u)(m')P(x,y),m''z)\\
  &=mf(u)(m')f(uv)(m'')P(P(x,y),z).
\end{aligned}
\]
Comparison of these equalities with \((*)\) gives
\(a(mx,m'y,m''z)=a(x,y,z)\).}
Condition \((+)_{3}\) therefore gives:

\smallskip
\noindent\textit{(1) There exists a unique morphism
\(DP:Y^3\to M\) such that}
\[
 P(x,P(y,z))=DP(p(x),p(y),p(z))\cdot P(P(x,y),z),
\]
\textit{and \(P\) is associative if and only if \(DP=0\).}

\smallskip
\noindent\textbf{(b)}
Let us compute \(P(P(P(x,y),z),t)\) step by step using the preceding
formula. Put
\[
 p(x)=u,\qquad p(y)=v,\qquad p(z)=w,\qquad p(t)=h.
\]
One obtains\footnote{Editor's note: the diagram that follows has been
added.} the following pentagonal diagram, in which an arrow
\(a\xrightarrow{m}b\) means that \(b=m\cdot a\).

\SGAThreeNativeDiagram[0.82\textwidth]
  {native_diagrams/exp3/diagram_030_exp3_p030_associativity_pentagon.tex}
 {The pentagonal comparison used to compute the associativity
  obstruction.}
 {fig:exp3-p030-associativity-pentagon}
 {LNM 151, pp.~109--110}

Thus one finds
\[
\begin{aligned}
 &DP(u,v,w)\cdot DP(u,vw,h)\cdot f(u)(DP(v,w,h))\\
 &\qquad\cdot DP(u,v,wh)^{-1}\cdot DP(uv,w,h)^{-1}=e_M,
\end{aligned}
\]
that is, \(\partial^3DP(u,v,w,h)=e_M\). Moreover, the left-hand side
of the preceding formula can be expressed, by means of \(P\) and
\(a\), as the value at \((x,y,z,t)\) of a certain morphism
\(X^4\to M\).

It follows from the uniqueness assertion in \((+)_{4}\) that
\(\partial^3DP\) and \(e_M\), which factor the same morphism, are
equal. Hence:

\smallskip
\noindent\textit{(2) \(DP\) is a cocycle; that is,
\(DP\in Z^3(Y,M)\).}

\smallskip
\noindent\textbf{(c)}
If \(P\) and \(P'\) are two admissible composition laws on \(X\),
there is a unique morphism
\[
 b:X^2\longrightarrow M
\]
such that \(P'(x,y)=b(x,y)P(x,y)\). Applying \SGARefLink{sga3:iii:definition:1-3-2:item:ii}{1.3.2(ii)} to \(P\) and
\(P'\), one sees that \(b\) is invariant under \(M^2\), and hence,
by \((+)_{2}\):

\smallskip
\noindent\textit{\SGARefLink{sga3:iii:section:3}{(3)} For every pair \((P,P')\) of admissible
composition laws there exists a unique morphism
\(d(P,P'):Y^2\to M\) such that}
\[
 P'(x,y)=d(P,P')(p(x),p(y))\cdot P(x,y),
\]
\textit{and the set of admissible composition laws is thereby made
principal homogeneous under}
\[
 \operatorname{Hom}(Y^2,M)=C^2(Y,M).
\]

\smallskip
\noindent\textbf{(d)}
Under the preceding conditions one has
\[
 (4)\qquad DP'-DP=\partial^2d(P,P').
\]

\smallskip
\noindent\textbf{\SGARefLink{sga3:iii:paragraph:2-0:equation:e}{(e)}}
The laws \(P\) and \(P'\) are equivalent if and only if there is a
morphism
\[
 h\in C^1(Y,M)=\operatorname{Hom}(Y,M)
\]
such that \(d(P,P')=\partial^1h\). This follows from the definition of
equivalence and from \SGARefLink{sga3:iii:lemma:1-3-4:item:ii}{1.3.4(ii)}.

\smallskip
\noindent\textbf{(f)}
It remains only to conclude. We seek an associative \(P'\), that is,
one for which \(DP'=e_M\). Now \(DP\) is a cocycle whose class in
\(H^3(Y,M)\) does not depend on the chosen admissible composition law
\(P\), by \SGARefLink{sga3:iii:section:3}{(3)} and \SGARefLink{sga3:iii:section:4}{(4)}. This class is the required obstruction \(c\).
An appropriate \(P'\) can be chosen if and only if \(c=0\). Indeed,
after choosing an arbitrary \(P\), one must solve, by \SGARefLink{sga3:iii:section:1}{(1)},
\[
 0=DP'=DP+\partial^2d(P,P'),
\]
which is possible, by \SGARefLink{sga3:iii:section:3}{(3)} and \SGARefLink{sga3:iii:section:4}{(4)}, if and only if \(c=0\). Again by
\SGARefLink{sga3:iii:section:3}{(3)} and \SGARefLink{sga3:iii:section:4}{(4)}, the set of associative laws \(P'\) is principal
homogeneous under \(Z^2(Y,M)\). By \SGARefLink{sga3:iii:paragraph:2-0:equation:e}{(e)}, the set of associative laws
\(P'\) modulo equivalence is principal homogeneous under
\(H^2(Y,M)\).

